\chapter{Aprendizaje en Entornos Complejos: De Métodos Tabulares a Aproximaciones Profundas}

\section{Introducción}

\subsection{Descripción del problema y su relevancia}
En esta segunda parte, nos centramos en el estudio de entornos secuenciales, 
donde la toma de decisiones no se limita a una única acción, sino que se extiende a lo largo del tiempo. 
Estos entornos, formalizados como Procesos de Decisión de Markov (MDP), representan un desafío significativo
en el campo del Aprendizaje por Refuerzo (RL) debido a la necesidad de planificar a largo 
plazo y manejar espacios de estados potencialmente infinitos.

Este tipo de técnicas tienen una gran relevancia debido a su aplicabilidad
en una amplia gama de dominios, desde la robótica y la conducción autónoma,
hasta la optimización de sistemas complejos y la toma de decisiones en finanzas.

\subsection{Motivación del trabajo}
La transición del problema de los bandidos a los Procesos de Decisión de Markov (MDP) 
introduce el reto de las recompensas retardadas y la planificación a largo plazo. 
Además, a medida que los problemas del mundo real se vuelven más intrincados, 
los espacios de estados crecen de manera exponencial. 
Esto hace que los enfoques clásicos resulten insuficientes debido a la limitación de 
memoria y tiempo de cómputo (la conocida ``maldición de la dimensionalidad''). 
Nuestra motivación radica en documentar y comprender empíricamente la transición
 metodológica necesaria para saltar de un entorno puramente discreto y manejable, 
 como \textit{Taxi-v3}, a un entorno con dinámicas físicas y espacios continuos, 
 como \textit{LunarLander}.

\subsection{Objetivos del capítulo}
El objetivo principal de este capítulo es implementar, evaluar y comparar diferentes algoritmos de Aprendizaje por Refuerzo capaces de manejar múltiples estados. De forma específica, buscamos:
\begin{itemize}
    \item Evaluar el rendimiento y la velocidad de convergencia de métodos de Monte Carlo y Diferencias Temporales en un espacio discreto.
    \item Implementar técnicas de aproximación de funciones para superar las limitaciones tabulares en espacios continuos.
    \item Contrastar empíricamente las propiedades \textit{on-policy} y \textit{off-policy} a través de distintas familias de algoritmos.
\end{itemize}

\subsection{Organización del capítulo}
El presente capítulo se estructura de la siguiente manera: la Sección 1 engloba esta introducción. La Sección 2 detalla el desarrollo y la formalización matemática del problema, junto con la literatura existente. La Sección 3 está dedicada a los métodos tabulares aplicados al entorno \textit{Taxi-v3}, abarcando Monte Carlo y Diferencias Temporales (SARSA, Expected-SARSA, Q-Learning y Double Q-Learning). Finalmente, la Sección 4 explora el control con aproximaciones en el entorno \textit{LunarLander}, donde implementamos SARSA semi-gradiente y Deep Q-Learning (DQN).

\section{Desarrollo}

\subsection{Definición formal del problema}
Para extender las ideas del bandido de $k$-brazos a situaciones secuenciales, formalizamos el problema de toma de decisiones como un Proceso de Decisión de Markov (MDP), definido por la tupla $\langle \mathcal{S}, \mathcal{A}, \mathcal{P}, \mathcal{R}, \gamma \rangle$. A diferencia del capítulo anterior, en cada instante de tiempo $t$, el agente observa un estado $S_t \in \mathcal{S}$ y ejecuta una acción $A_t \in \mathcal{A}$. Como consecuencia, el entorno transita a un nuevo estado $S_{t+1}$ según la probabilidad de transición $\mathcal{P}(S_{t+1} | S_t, A_t)$ y emite una recompensa escalar $R_{t+1} \in \mathcal{R}$. Nuestro objetivo es encontrar una política $\pi(a|s)$ que maximice el retorno esperado descontado:
\begin{equation}
    G_t = \sum_{k=0}^{\infty} \gamma^k R_{t+k+1}
\end{equation}
donde $\gamma \in [0, 1]$ es el factor de descuento que pondera la importancia de las recompensas futuras.

% 

\subsection{Enfoques existentes en la literatura}
La literatura clásica de RL establece dos grandes ramas para la resolución de MDPs: los métodos tabulares y los métodos de aproximación. Para espacios discretos pequeños, algoritmos como Q-Learning y SARSA garantizan la convergencia a la política óptima construyendo un diccionario exacto de valores. Sin embargo, para entornos continuos o de alta dimensionalidad, la literatura reciente se apoya en aproximadores de funciones no lineales. La introducción de Deep Q-Networks (DQN) por Mnih et al. marcó un punto de inflexión al demostrar que las redes neuronales profundas pueden estabilizar el aprendizaje \textit{off-policy} mediante el uso de repetición de experiencia (\textit{Experience Replay}) y redes objetivo (\textit{Target Networks}).

\subsection{Contexto y antecedentes necesarios}
Para comprender la progresión de nuestro trabajo en este capítulo, es vital establecer las características de los dos entornos de prueba que utilizaremos para superar el nivel del bandido clásico:
\begin{itemize}
    \item \textbf{Taxi-v3:} Un entorno discreto clásico con 500 estados posibles y 6 acciones discretas. Permite el uso de matrices exactas ($Q(s, a)$) para estimar el valor de cada par estado-acción.
    \item \textbf{LunarLander:} Un simulador físico donde el estado es un vector continuo de 8 dimensiones (coordenadas, velocidades, ángulos, etc.). Aquí, almacenar una tabla es imposible, lo que nos obliga a estimar la función de valor mediante pesos parametrizados: $\hat{q}(s, a, \mathbf{w}) \approx q_\pi(s, a)$.
\end{itemize}

\subsection{Explicación de los métodos utilizados}
Para abordar estos retos, dividimos nuestra metodología en dos bloques principales:
\begin{enumerate}
    \item \textbf{Métodos Tabulares:} Aplicamos Monte Carlo para aprender de episodios completos y Diferencias Temporales (TD) para aprender paso a paso sin esperar al final del episodio. Exploramos el control \textit{on-policy} (SARSA, Expected-SARSA) que optimiza la política de comportamiento, y el control \textit{off-policy} (Q-Learning, Double Q-Learning) que busca la política óptima independientemente de la exploración.
    \item \textbf{Aproximaciones:} Reemplazamos la tabla $Q$ por un modelo parametrizado. Iniciamos con SARSA semi-gradiente utilizando características construidas manualmente o combinaciones lineales, para finalmente evolucionar a Deep Q-Learning (DQN), delegando la extracción de características a una red neuronal profunda.
\end{enumerate}

\subsection{Justificación y diferenciación del trabajo}
A diferencia de estudios teóricos aislados, el presente capítulo destaca por su enfoque integrador y comparativo. Justificamos la estructura al demostrar empíricamente cómo las limitaciones metodológicas de los enfoques tabulares (analizados en \textit{Taxi-v3}) se resuelven mediante el uso de arquitecturas neuronales en \textit{LunarLander}. Al implementar las variaciones algorítmicas bajo un marco de pruebas unificado, aportamos una visión clara y cuantificable sobre el equilibrio entre la simplicidad de los métodos clásicos y la capacidad de generalización de los métodos profundos contemporáneos.